% Part: incompleteness
% Chapter: incompleteness-provability
% Section: fixed-point-lemma

\documentclass[../../../include/open-logic-section]{subfiles}

\begin{document}

\olfileid{inc}{inp}{fix}

\olsection{د ثابت ټکي لمه}

\begin{explain}
د ثابت ټکي لمه وايي چې د هر !!{formula}~$!B(x)$ دپاره داسې
!!a{sentence}~$!A$ شته چې
$\Th{T} \Proves !A \liff !B(\gn{!A})$، په دې شرط چې~$\Th{T}$
تيوري~$\Th{Q}$ غځوي۔ د دروغجنې جملې په بېلګه کښې غواړو چې
$!A$ په~$\Th{T}$ کښې په ثابتېدونکي ډول له «$\gn{!A}$ ناسمه ده»
سره معادله وي؛ يعنې له هغې جملې سره چې وايي~$\Gn{!A}$ د يوې
ناسمې !!{sentence} ګوډل شمېره ده۔ د ثبوت د مفکورې د پوهېدو دپاره
دا ګټوره ده چې هغه د~$!A$ په باب د کواین له دې غير صوري څرګندونې
سره پرتله کړو:
«که يې خپله نقل شوې بڼه ترې مخکې کېښودل شي، دروغ وايي» که يې
خپله نقل شوې بڼه ترې مخکې کېښودل شي، دروغ وايي۔ له يوه عبارت څخه
د هغه خپله نقل شوې بڼه مخکې ايښودل او بيا ترې جمله جوړول د هغه
عبارت \emph{قطرول} بللے شو، او پايله ئې د هغه قطري بڼه ده۔ نو د
«که يې خپله نقل شوې بڼه ترې مخکې کېښودل شي، دروغ وايي» قطري بڼه
همدا «که يې خپله نقل شوې بڼه ترې مخکې کېښودل شي، دروغ وايي» که يې
خپله نقل شوې بڼه ترې مخکې کېښودل شي، دروغ وايي، ده۔ اوس پام وکړئ
چې د کواین دروغجنه جمله د «دروغ وايي» قطري بڼه نۀ ده، بلکې د
«که يې خپله نقل شوې بڼه ترې مخکې کېښودل شي، دروغ وايي» قطري بڼه
ده۔ نو هغه خاصيت چې د دروغجنې جملې د جوړولو دپاره قطرول کېږي،
پخپله هم قطرول رانغاړي!{}

د حساب په ژبه کښې د يو آزاد متغير لرونکي !!a{formula} نقل شوې بڼه
د هغه د ګوډل شمېرې په حسابولو او بيا د دغې شمېرې د معياري عدديې
په آزاد متغير کښې په ځاے کولو جوړوو۔ د~$!E(x)$ قطري بڼه
$!E(\num{n})$ ده، چېرته چې $n = \Gn{!E(x)}$ دے۔ (له دې وروسته
$\num{\Gn{!E(x)}}$ په~$\gn{!E(x)}$ لنډوو۔) نو که $!B(x)$ د
«ناسمه ده» خاصيت وي، «که خپله نقل شوې بڼه ترې مخکې راشي، دروغ وايي»
به دا وي: «کله چې د خپلې قطري بڼې پر ګوډل شمېرې پلے شي، دروغ وايي»۔
که د هغې تابع~$\fn{diag}(n)$ دپاره چې د ګوډل شمېرې~$n$ لرونکي
!!{formula} د قطري بڼې ګوډل شمېره حسابوي، يوه نښه~$\Obj{diag}$
ولرو، $!E(x)$ به د~$!B(\Obj{diag}(x))$ په بڼه وليکو۔ د کواین
د دروغجنې جملې بڼه به بيا د همدې عبارت قطري بڼه وي، يعنې
$!E(\gn{!E(x)})$ يا
$!B(\Obj{diag}(\gn{!B(\Obj{diag}(x))}))$۔ البته اوس $!B(x)$
هر بل خاصيت هم کېدے شي، او همدا جوړښت به بيا هم کار وکړي۔ د
ناتکميلۍ د قضيې دپاره $!B(x)$ داسې اخلو: «$x$ په~$\Th{T}$ کښې
!!{derivable} نۀ دے»۔ نو $!E(x)$ به ووايي: «کله چې د خپلې قطري
بڼې پر ګوډل شمېرې پلے شي، داسې !!a{sentence} جوړوي چې په~$\Th{T}$
کښې !!{derivable} نۀ ده»۔

د دې خبرې په~$\Th{T}$ کښې د صوري کولو دپاره بايد د~$\fn{diag}$
د صوري کولو لار ومومو۔ تابع~$\fn{diag}(n)$ محاسبه کېدونکې ده، بلکې
ابتدايي بازګشتي ده: که $n$ د يوه فورمول~$!E(x)$ ګوډل شمېره وي،
$\fn{diag}(n)$ د~$!E(\gn{!E(x)})$ ګوډل شمېره ورکوي۔ (ياد ولرئ
چې~$\gn{!E(x)}$ د~$!E(x)$ د ګوډل شمېرې معياري عدديه، يعنې
$\num{\Gn{!E(x)}}$، ده۔) که $\Obj{diag}$ په~$\Th{T}$ کښې د
تابع~$\fn{diag}$ تمثيلوونکې تابع‌نښه واے، موږ به $!A$ د
$!B(\Obj{diag}(\gn{!B(\Obj{diag}(x))}))$ فورمول اخيستے واے۔
پام وکړئ چې
\begin{align*}
\fn{diag}(\Gn{!B(\Obj{diag}(x))}) & =
\Gn{!B(\Obj{diag}(\gn{!B(\Obj{diag}(x))}))} \\
& = \Gn{!A}.
\end{align*}
که فرض کړو چې~$\Th{T}$
\[
\Obj{diag}(\gn{!B(\Obj{diag}(x))}) = \gn{!A},
\]
!!{derive} کولے شي، نو دا هم !!{derive} کولے شي چې
$!B(\Obj{diag}(\gn{!B(\Obj{diag}(x))}))
\liff !B(\gn{!A})$۔ خو د دې معادلۍ کيڼ اړخ د تعريف له مخې~$!A$
دے۔

البته $\Obj{diag}$ په عمومي توګه د~$\Th{T}$ تابع‌نښه نۀ ده، او
يقيناً د~$\Th{Q}$ تابع‌نښه نۀ ده۔ خو ځکه چې~$\fn{diag}$ محاسبه
کېدونکې ده، په~$\Th{Q}$ کښې د کوم فورمول
$!D_{\fn{diag}}(x,y)$ په وسيله \emph{تمثيلېږي}۔ نو د
$!B(\Obj{diag}(x))$ پر ځاے
$\lexists[y][(!D_{\fn{diag}}(x,y) \land !B(y))]$ ليکلے شو۔
له دې پرته، پورته لنډ بيان شوے ثبوت بشپړېږي، او په حقيقت کښې
لا په~$\Th{Q}$ کښې بشپړېږي۔
\end{explain}

\begin{lem}
\ollabel{lem:fixed-point} فرض کړئ $!B(x)$ هر هغه فورمول دے چې
يو آزاد متغير~$x$ لري۔ نو داسې جمله~$!A$ شته چې
$\Th{Q} \Proves !A \liff !B(\gn{!A})$۔
\end{lem}

\begin{proof}
د ورکړي~$!B(x)$ دپاره دې $!E(x)$ فورمول
$\lexists[y][(!D_{\fn{diag}}(x,y) \land !B(y))]$ وي، او
$!A$ دې د هغه قطري بڼه، يعنې $!E(\gn{!E(x)})$ فورمول وي۔

ځکه چې $!D_{\fn{diag}}$ د~$\fn{diag}$ تمثيل کوي، او
$\fn{diag}(\Gn{!E(x)}) = \Gn{!A}$، نو~$\Th{Q}$ دا
!!{derive} کولے شي:
\begin{align}
  & !D_{\fn{diag}}(\gn{!E(x)}, \gn{!A}) \ollabel{repdiag1} \\
  & \lforall[y][(!D_{\fn{diag}}(\gn{!E(x)},y) \lif
  \eq[y][\gn{!A}])]. \ollabel{repdiag2}
\end{align}
اوس ښيو چې $\Th{Q} \Proves !A \liff !B(\gn{!A})$۔ موږ يوازې د
منطق او په~$\Th{Q}$ کښې د !!{derivable} حقايقو په کارولو، په غير
صوري ډول استدلال کوو۔

لومړے فرض کړئ~$!A$، يعنې $!E(\gn{!E(x)})$۔ د~$!E(x)$ تعريف
ته په ستنېدو وينو چې $!E(\gn{!E(x)})$ همدا دے:
\[
\lexists[y][(!D_{\fn{diag}}(\gn{!E(x)},y) \land !B(y))].
\]
داسې يو~$y$ وګورئ۔ ځکه چې
$!D_{\fn{diag}}(\gn{!E(x)},y)$ دے، د~\olref{repdiag2} له مخې
$y = \gn{!A}$۔ نو له~$!B(y)$ څخه~$!B(\gn{!A})$ اخلو۔

اوس فرض کړئ $!B(\gn{!A})$۔ د~\olref{repdiag1} له مخې لرو:
\begin{align*}
& !D_{\fn{diag}}(\gn{!E(x)}, \gn{!A}) \land !B(\gn{!A}).
\intertext{نو ورڅخه راځي چې}
& \lexists[y][(!D_{\fn{diag}}(\gn{!E(x)},y) \land !B(y))].
\end{align*}
خو دا هماغه $!E(\gn{!E(x)})$، يعنې~$!A$، دے۔
\end{proof}

\begin{digress}
دا ثبوت د محاسبه کېدنې په تيورۍ کښې د ثابت ټکي د لمې له ثبوت سره
پرتله کړئ۔ توپير دا دے چې دلته غواړو يو \emph{بيان} د خپل ځان په
وسيله تعريف کړو، خو هلته مو يوه \emph{تابعه} د خپل ځان په وسيله
تعريفوله؛ له دې توپيره پرته مفکوره په رښتيا هماغه ده۔
\end{digress}

\begin{prob}
  !!^a{formula}~$!A(x)$ د \emph{رښتياوالي تعريف} دے، که د هرې
  !!{sentence}s~$!B$ دپاره $\Th{Q} \Proves !B \liff !A(\gn{!B})$ وي۔
  د ثابت ټکي د لمې په کارولو وښيئ چې هېڅ !!{formula} د رښتياوالي
  تعريف نۀ دے۔
\end{prob}

\end{document}
